%% ================================================================
%%  Preamble -- target: Nature Computational Science (Article)
%%  Limits (third-party guide, verify on nature.com):
%%    main text <= 3,500 words (excl. abstract, Methods, refs, legends)
%%    abstract <= 150 words, unreferenced
%%    <= 6 display items (figures + tables)
%% ================================================================

%% ---------- Packages ----------
\usepackage[T1]{fontenc}
\usepackage[utf8]{inputenc}
\usepackage{lmodern}
\usepackage[a4paper,margin=2.5cm]{geometry}
\usepackage{setspace}
\usepackage{amsmath, amssymb, amsthm}
\usepackage{mathtools}
\usepackage{bm}
\usepackage{graphicx}
\usepackage{subcaption}
\usepackage{booktabs}
\usepackage{xcolor}
\usepackage{algorithm}
\usepackage{algpseudocode}
\usepackage[numbers,sort&compress]{natbib}
\usepackage[colorlinks=true,linkcolor=blue!60!black,citecolor=blue!60!black,urlcolor=blue!60!black]{hyperref}
\usepackage{cleveref}   % after hyperref
\usepackage{authblk}

\graphicspath{{figures/}}

%% ---------- Notation (same as MAIN.tex and the workshop draft) ----------
\newcommand{\zstar}{z^*}
\newcommand{\ux}{u^*}
\newcommand{\uhat}{\hat{u}}
\newcommand{\bw}{\mathbf{w}}
\newcommand{\bu}{\mathbf{u}}
\newcommand{\bA}{\mathbf{A}}
\newcommand{\bb}{\mathbf{b}}
\newcommand{\R}{\mathbb{R}}
\newcommand{\Lop}{\mathcal{L}}
\newcommand{\Mset}{\mathcal{M}}          % the map (set of points)
\newcommand{\Sset}{\mathcal{S}}          % a stencil
\newcommand{\Aset}{\mathcal{A}}          % action set (menu)
\newcommand{\normone}[1]{\lVert #1 \rVert_1}
\newcommand{\normtwo}[1]{\lVert #1 \rVert_2}
\newcommand{\EPS}{\mathrm{EPS}}
\newcommand{\Nref}{N_{\mathrm{ref}}}

%% ---------- Drafting markers ----------
%% \NUM{}  = number taken from a notebook/run file; re-check before submitting
%% \CITE{} = citation not yet resolved to a bibtex key
%% \TODO{} = missing content; stays red in final mode on purpose
\newif\ifdraftmode
\draftmodetrue          % <<-- \draftmodefalse for the submission build

\newcommand{\TODO}[1]{\textcolor{red}{\textbf{[TODO: #1]}}}
\newcommand{\CITE}[1]{\textcolor{orange}{\textbf{[cite: #1]}}}
\ifdraftmode
  \newcommand{\NUM}[1]{\textcolor{blue}{\underline{#1}}}
\else
  \newcommand{\NUM}[1]{#1}
\fi

%% ---------- Theorem environments ----------
\theoremstyle{plain}
\newtheorem{proposition}{Proposition}
\theoremstyle{remark}
\newtheorem{remark}{Remark}
